% ============================================================
\chapter{Dualit\'e de Fenchel--Moreau--Rockafellar}
\label{ch:fenchel}
% ============================================================

La transform\'ee de Legendre, connue des physiciens depuis le XVIII\textsuperscript{e} si\`ecle pour passer du lagrangien au hamiltonien, a \'et\'e g\'en\'eralis\'ee par Werner Fenchel dans les ann\'ees 1940 en un outil d'une puissance remarquable~: la \emph{conjugaison convexe}. L'id\'ee est de repr\'esenter une fonction convexe non pas par son graphe, mais par l'ensemble de ses hyperplans d'appui --- une sorte de dualit\'e g\'eom\'etrique o\`u chaque fonction convexe a un alter ego. Le th\'eor\`eme de biconjugaison de Fenchel-Moreau affirme que pour une fonction convexe semi-continue inf\'erieurement, appliquer deux fois cette transformation redonne la fonction de d\'epart. Cette involution est la cl\'e de la dualit\'e en optimisation convexe.

\section{Introduction}

La conjugaison convexe, ou transform\'ee de Legendre--Fenchel, est un outil fondamental
de l'analyse convexe. Elle \'etablit une correspondance entre fonctions convexes et
permet de construire des probl\`emes duaux puissants. Le th\'eor\`eme de biconjugaison
de Fenchel--Moreau caract\'erise les fonctions convexes semi-continues inf\'erieurement.

\section{Conjugu\'ee de Fenchel}

\begin{definition}[Conjugu\'ee de Fenchel]
Soit $f : \R^n \to \R \cup \{+\infty\}$. La \emph{conjugu\'ee de Fenchel} (ou
\emph{transform\'ee de Legendre--Fenchel}) de $f$ est~:
\[
    f^*(y) = \sup_{x \in \R^n} \left\{\inner{y}{x} - f(x)\right\}.
\]
\end{definition}

\begin{proposition}[Propri\'et\'es imm\'ediates]
\begin{enumerate}
    \item $f^*$ est toujours convexe et semi-continue inf\'erieurement (comme supremum
          de fonctions affines en $y$).
    \item \textbf{In\'egalit\'e de Fenchel--Young}~:
    $f(x) + f^*(y) \geq \inner{x}{y}$ pour tout $x, y$.
    \item Si $f$ est propre, alors $f^*$ ne prend jamais la valeur $-\infty$.
\end{enumerate}
\end{proposition}

\begin{proof}[Preuve de l'in\'egalit\'e de Fenchel--Young]
Par d\'efinition de $f^*$~:
\[
    f^*(y) = \sup_{z} \{\inner{y}{z} - f(z)\} \geq \inner{y}{x} - f(x),
\]
d'o\`u $f(x) + f^*(y) \geq \inner{x}{y}$.
\end{proof}

\begin{formules}[title=\textbf{In\'egalit\'e de Fenchel--Young}]
\[
    f(x) + f^*(y) \geq \inner{x}{y}, \quad \forall x, y \in \R^n,
\]
avec \'egalit\'e si et seulement si $y \in \partial f(x)$ (ou de mani\`ere \'equivalente
$x \in \partial f^*(y)$).
\end{formules}

\section{Exemples de conjugu\'ees}

\begin{exemple}[Conjugu\'ees classiques]
\begin{enumerate}
    \item \textbf{Fonction affine}~: $f(x) = \inner{a}{x} + b$.
    Alors $f^*(y) = \begin{cases} -b & \text{si } y = a, \\ +\infty & \text{sinon.}\end{cases}$

    \item \textbf{Quadratique}~: $f(x) = \frac{1}{2}x^TQx$ avec $Q \succ 0$.
    Alors $f^*(y) = \frac{1}{2}y^TQ^{-1}y$.

    \item \textbf{Norme}~: $f(x) = \norm{x}_p$ pour $p \geq 1$.
    Alors $f^* = \iota_{B_q}$ o\`u $B_q = \{y \mid \norm{y}_q \leq 1\}$ et
    $1/p + 1/q = 1$.

    \item \textbf{Norme au carr\'e}~: $f(x) = \frac{1}{2}\norm{x}^2$.
    Alors $f^*(y) = \frac{1}{2}\norm{y}^2$ (auto-conjugu\'ee).

    \item \textbf{Indicatrice}~: $f = \iota_C$.
    Alors $f^* = \sigma_C$ (fonction d'appui).

    \item \textbf{Entropie n\'egative}~: $f(x) = x\log x - x$ pour $x > 0$.
    Alors $f^*(y) = e^y$.

    \item \textbf{Log-sum-exp}~: $f(x) = \log(\sum_i e^{x_i})$.
    Alors $f^*(y) = \sum_i y_i \log y_i$ si $y \in \Delta_n$, $+\infty$ sinon.
\end{enumerate}
\end{exemple}

\begin{proof}[Preuve pour la quadratique]
$f^*(y) = \sup_x \{\inner{y}{x} - \frac{1}{2}x^TQx\}$. La condition d'optimalit\'e
donne $y - Qx = 0$, soit $x = Q^{-1}y$. Donc~:
\[
    f^*(y) = \inner{y}{Q^{-1}y} - \frac{1}{2}(Q^{-1}y)^TQ(Q^{-1}y)
    = y^TQ^{-1}y - \frac{1}{2}y^TQ^{-1}y = \frac{1}{2}y^TQ^{-1}y.
\]
\end{proof}

\section{Th\'eor\`eme de biconjugaison de Fenchel--Moreau}

\begin{theoreme}[Fenchel--Moreau]
\label{thm:fenchel_moreau}
Soit $f : \R^n \to \R \cup \{+\infty\}$ propre. Alors~:
\[
    f^{**} = f \quad \Longleftrightarrow \quad
    f \text{ est convexe et semi-continue inf\'erieurement.}
\]
En g\'en\'eral, $f^{**}$ est la plus grande fonction convexe s.c.i.\ minor\'ee par $f$
(la \emph{r\'egularis\'ee convexe s.c.i.} de $f$, not\'ee $\overline{\mathrm{conv}}(f)$).
\end{theoreme}

\begin{proof}[Id\'ee de la preuve]
On montre d'abord que $f^{**} \leq f$ (par l'in\'egalit\'e de Fenchel--Young).
Ensuite, on utilise le th\'eor\`eme de s\'eparation~: si $f$ est convexe s.c.i.\
et $(x_0, t_0) \notin \mathrm{epi}(f)$, il existe un hyperplan s\'eparant
$(x_0, t_0)$ de $\mathrm{epi}(f)$, ce qui montre $f^{**}(x_0) \geq t_0$.
Comme ceci vaut pour tout $t_0 < f(x_0)$, on obtient $f^{**} \geq f$.
\end{proof}

\section{Lien entre conjugu\'ee et sous-diff\'erentiel}

\begin{theoreme}[\'Equivalence conjugu\'ee--sous-diff\'erentiel]
\label{thm:conjuguee_subdiff}
Soit $f$ convexe propre et s.c.i. Les assertions suivantes sont \'equivalentes~:
\begin{enumerate}
    \item $y \in \partial f(x)$,
    \item $x \in \partial f^*(y)$,
    \item $f(x) + f^*(y) = \inner{x}{y}$.
\end{enumerate}
\end{theoreme}

\begin{proof}
$(1) \Rightarrow (3)$~: Si $y \in \partial f(x)$, alors pour tout $z$~:
$f(z) \geq f(x) + \inner{y}{z-x}$, soit $\inner{y}{z} - f(z) \leq \inner{y}{x} - f(x)$.
Donc $f^*(y) = \inner{y}{x} - f(x)$.

$(3) \Rightarrow (1)$~: $f^*(y) = \inner{y}{x} - f(x)$ signifie que le supremum dans
la d\'efinition de $f^*$ est atteint en $x$, donc $y \in \partial f(x)$.

$(1) \Leftrightarrow (2)$~: Par $f^{**} = f$ et sym\'etrie.
\end{proof}

\section{Dualit\'e de Fenchel--Rockafellar}

\begin{theoreme}[Dualit\'e de Fenchel--Rockafellar]
\label{thm:fenchel_rockafellar}
Soient $f, g : \R^n \to \R \cup \{+\infty\}$ convexes propres et $A \in \R^{m \times n}$.
Consid\'erons le probl\`eme primal~:
\[
    (P) \quad p^* = \inf_{x \in \R^n} \{f(x) + g(Ax)\}
\]
et le probl\`eme dual~:
\[
    (D) \quad d^* = \sup_{y \in \R^m} \{-f^*(-A^Ty) - g^*(y)\}.
\]
Alors~:
\begin{enumerate}
    \item (Dualit\'e faible) $d^* \leq p^*$.
    \item (Dualit\'e forte) Si $\mathrm{ri}(\mathrm{dom}(f)) \cap
          A^{-1}(\mathrm{ri}(\mathrm{dom}(g))) \neq \emptyset$ (condition de
          qualification), alors $d^* = p^*$ et le supremum dual est atteint.
\end{enumerate}
\end{theoreme}

\begin{proof}[Preuve de la dualit\'e faible]
Pour tout $x$ et $y$, par l'in\'egalit\'e de Fenchel--Young~:
\begin{align*}
    f(x) &\geq \inner{-A^Ty}{x} - f^*(-A^Ty) = -\inner{y}{Ax} - f^*(-A^Ty), \\
    g(Ax) &\geq \inner{y}{Ax} - g^*(y).
\end{align*}
En sommant~: $f(x) + g(Ax) \geq -f^*(-A^Ty) - g^*(y)$.
En prenant l'infimum en $x$ et le supremum en $y$~: $p^* \geq d^*$.
\end{proof}

\section{Cas particulier~: $A = I$}

\begin{corollaire}[Dualit\'e de Fenchel]
Pour $f, g$ convexes propres avec $\mathrm{ri}(\mathrm{dom}(f)) \cap
\mathrm{ri}(\mathrm{dom}(g)) \neq \emptyset$~:
\[
    \inf_x \{f(x) + g(x)\} = \sup_y \{-f^*(-y) - g^*(y)\}
    = -\inf_y \{f^*(-y) + g^*(y)\}.
\]
\end{corollaire}

\section{R\`egles de calcul des conjugu\'ees}

\begin{proposition}[R\`egles de calcul]
\begin{enumerate}
    \item \textbf{S\'eparabilit\'e}~: si $f(x) = \sum_i f_i(x_i)$, alors
          $f^*(y) = \sum_i f_i^*(y_i)$.
    \item \textbf{Mise \`a l'\'echelle}~: $(\alpha f)^*(y) = \alpha f^*(y/\alpha)$
          pour $\alpha > 0$.
    \item \textbf{Translation}~: si $g(x) = f(x - a)$, alors $g^*(y) = f^*(y) + \inner{y}{a}$.
    \item \textbf{Ajout de lin\'eaire}~: si $g(x) = f(x) + \inner{b}{x} + c$,
          alors $g^*(y) = f^*(y - b) - c$.
    \item \textbf{Composition lin\'eaire}~: si $g(x) = f(Ax)$ avec $A$ inversible,
          alors $g^*(y) = f^*(A^{-T}y)$.
    \item \textbf{Infimale convolution}~: $(f \mathop{\square} g)^* = f^* + g^*$.
\end{enumerate}
\end{proposition}

\section{Application~: formulation duale du LASSO}

Le probl\`eme LASSO primal est~:
\[
    \min_x \frac{1}{2}\norm{Ax - b}^2 + \lambda\norm{x}_1.
\]

En \'ecrivant $f(x) = \lambda\norm{x}_1$ et $g(z) = \frac{1}{2}\norm{z - b}^2$
avec $z = Ax$~:

\begin{align*}
    f^*(y) &= \iota_{\{y \mid \norm{y}_\infty \leq \lambda\}}(y), \\
    g^*(w) &= \frac{1}{2}\norm{w}^2 + \inner{w}{b}.
\end{align*}

Le dual de Fenchel--Rockafellar donne~:
\[
    \max_\nu \left\{-\frac{1}{2}\norm{\nu}^2 + \inner{\nu}{b} - \frac{1}{2}\norm{b}^2\right\}
    \quad \text{sous} \quad \norm{A^T\nu}_\infty \leq \lambda.
\]

Soit, apr\`es simplification~:
\[
    \max_\nu \left\{-\frac{1}{2}\norm{b - \nu}^2\right\}
    \quad \text{sous} \quad \norm{A^T\nu}_\infty \leq \lambda.
\]

\section{Impl\'ementation Python}

\begin{codeblock}[title=\textbf{Calcul de conjugu\'ees et v\'erification de Fenchel--Young}]
\begin{minted}{python}
import numpy as np
import cvxpy as cp

def fenchel_conjugate_numerical(f, y, n, bounds=(-10, 10)):
    """Calcul numerique de f*(y) = sup_x {<y,x> - f(x)}."""
    x = cp.Variable(n)
    objective = cp.Maximize(y @ x - f(x))
    prob = cp.Problem(objective)
    prob.solve()
    return prob.value

# Exemple: f(x) = ||x||_2^2 / 2
n = 5
y_test = np.random.randn(n)

f = lambda x: 0.5 * cp.sum_squares(x)
f_star_val = fenchel_conjugate_numerical(f, y_test, n)
f_star_exact = 0.5 * np.sum(y_test**2)

print(f"f*(y) numerique:  {f_star_val:.6f}")
print(f"f*(y) analytique: {f_star_exact:.6f}")

# Verification Fenchel-Young: f(x) + f*(y) >= <x,y>
x_test = np.random.randn(n)
f_x = 0.5 * np.sum(x_test**2)
lhs = f_x + f_star_exact
rhs = np.dot(x_test, y_test)
print(f"f(x) + f*(y) = {lhs:.4f} >= <x,y> = {rhs:.4f}: {lhs >= rhs - 1e-10}")

# Dualite de Fenchel pour le LASSO
m, nn = 50, 20
A = np.random.randn(m, nn)
b = np.random.randn(m)
lam = 1.0

# Primal
x = cp.Variable(nn)
primal = cp.Problem(cp.Minimize(
    0.5 * cp.sum_squares(A @ x - b) + lam * cp.norm1(x)))
primal.solve()
print(f"\nLASSO primal: {primal.value:.6f}")

# Dual
nu = cp.Variable(m)
dual = cp.Problem(cp.Maximize(
    -0.5 * cp.sum_squares(b - nu)),
    [cp.norm_inf(A.T @ nu) <= lam])
dual.solve()
print(f"LASSO dual:   {dual.value + 0.5*np.sum(b**2):.6f}")
\end{minted}
\end{codeblock}

\section{Exercices}

\begin{exercice}[$\star$]
Calculer la conjugu\'ee de $f(x) = \frac{1}{p}\norm{x}_p^p$ pour $p > 1$.
V\'erifier que $f^*(y) = \frac{1}{q}\norm{y}_q^q$ avec $1/p + 1/q = 1$.
\end{exercice}

\begin{exercice}[$\star$]
Montrer que si $f(x) = \iota_C(x)$, alors $f^* = \sigma_C$.
\end{exercice}

\begin{exercice}[$\star\star$]
Montrer que $f^{**} \leq f$ toujours, et que $f^{**} = f$ si et seulement si $f$
est convexe et semi-continue inf\'erieurement.
\end{exercice}

\begin{exercice}[$\star\star$]
Calculer la conjugu\'ee de $f(x) = \log(1 + e^x)$ et interpr\'eter le r\'esultat
en termes d'entropie.
\end{exercice}

\begin{exercice}[$\star\star$]
D\'eriver le probl\`eme dual du LASSO par la dualit\'e de Fenchel--Rockafellar.
\end{exercice}

\begin{exercice}[$\star\star\star$]
D\'emontrer le th\'eor\`eme de Fenchel--Moreau en utilisant le th\'eor\`eme de
s\'eparation par hyperplan.
\end{exercice}

\begin{exercice}[$\star\star\star$]
Montrer la r\`egle $(f \mathop{\square} g)^* = f^* + g^*$ et en d\'eduire que
l'infimale convolution de deux fonctions convexes propres s.c.i.\ est convexe s.c.i.
\end{exercice}
